% Vietnamese translation for OI Public Library
% Source-File: IOI中国国家候选队论文/2009/姜碧野/SPFA算法的优化及应用.ppt
\documentclass[11pt]{article}
\usepackage{oipl-vn}

\newcommand{\Oh}{\mathcal{O}}

\title{Công cụ mạnh của lời giải lặp: tối ưu hóa và ứng dụng thuật toán SPFA\\{\large Ghi chú theo slide}}
\author{Jiang Biye\\Trường Trung học Kỷ niệm Tôn Trung Sơn, Quảng Đông}
\date{Luận văn Trại đông 2009}

\begin{document}
\maketitle

\origtitle{Optimization and applications of SPFA}
\sourcepath{IOI-China-Candidate-Team-Papers/2009/Jiang-Biye/spfa-slides.ppt}

\section{Mạch trình bày}

Slide xem SPFA không chỉ là thuật toán đường đi ngắn nhất, mà còn là khuôn
mẫu cho các quá trình nới lỏng lặp. Vì vậy nội dung được chia thành ba tầng:
cài đặt cơ bản, tối ưu thực tế, và ứng dụng vào các mô hình không hiển nhiên
là đồ thị.

\section{Các điểm cần giữ}

\begin{itemize}
  \item SPFA duy trì hàng đợi các đỉnh có khả năng làm thay đổi nhãn của đỉnh
  khác. Mỗi lần lấy một đỉnh ra, ta thử nới lỏng các cạnh kề.
  \item Cài đặt DFS có thể phát hiện chu trình âm hoặc dương nhanh trong một
  số hệ ràng buộc, vì chu trình được gặp ngay trên ngăn đệ quy.
  \item Các tối ưu như SLF, LLL, đánh dấu trạng thái trong hàng đợi và bỏ qua
  trạng thái vô dụng đều nhằm giảm số lần nới lỏng thừa.
  \item Trong hệ ràng buộc hiệu, mỗi bất đẳng thức được đổi thành một cạnh;
  nghiệm khả thi tương ứng với một bộ nhãn thỏa toàn bộ cạnh.
\end{itemize}

\section{Ứng dụng được nhấn mạnh}

Các bài quy hoạch động có thứ tự chuyển trạng thái khó xác định có thể được
xem như bài toán nới lỏng trên đồ thị trạng thái. Khi một trạng thái tốt hơn,
những trạng thái phụ thuộc vào nó được đưa lại vào hàng đợi. Cách nhìn này
giúp thay thế thứ tự duyệt cứng bằng một cơ chế lan truyền thay đổi.

\section{Khi đọc cùng bài viết}

Bài viết đầy đủ có các đề ví dụ và dữ liệu thử. Bản slide nên được dùng để
nắm thông điệp chính: SPFA hữu ích khi ta có quan hệ nới lỏng cục bộ, số thay
đổi thực tế ít hơn nhiều so với số cạnh lý thuyết, và việc phát hiện sớm
chu trình hoặc trạng thái vô dụng làm thay đổi đáng kể thời gian chạy.

\end{document}
